% ============================================================
% Section 3: Method
% Draft 1. Every number traced to a results document in the repo.
% Normative references: docs/RESIDUAL.md (residual-v3) and
% docs/SPEC-PERTURBATION-MODEL.md (perturbation-v1) only.
% ============================================================

\section{Method}
\label{sec:method}

This section defines four things: the residual statistic under test, the
perturbation model under which it was exercised, the threshold against which
detection was scored, and the discipline used to decide what counts as a
result. The first three are ordinary experimental description. The fourth
matters more than usual here, because most of what follows is a negative
result, and a negative result is only worth reading if the bar it failed was
set before the data arrived.

\subsection{The residual statistic}
\label{sec:method:residual}

The detector is a one-sided checksum in the Huang--Abraham family. Let
$A \in \mathbb{R}^{M \times K}$ and $B \in \mathbb{R}^{K \times N}$ be the
factors, $C$ the computed product, and $e = \mathbf{1}_N$. Two length-$M$
vectors are formed,
%
\begin{equation}
d = C e, \qquad d' = A (B e),
\end{equation}
%
which agree exactly in exact arithmetic. The detector compares them per row,
normalized by a nonnegative contraction of the factors:
%
\begin{equation}
\label{eq:residual-v3}
r_i = \frac{\lvert d_i - d'_i \rvert}{s_i},
\qquad
s = \lvert A \rvert \bigl( \lvert B \rvert e \bigr),
\qquad
r_{\max} = \max_i r_i .
\end{equation}
%
Absolute values are entrywise. All of $A$, $B$ and $C$ are promoted to
\texttt{float64} before these operations, and every accumulation in
Equation~\ref{eq:residual-v3} is performed in \texttt{float64}. We refer to
this statistic as \emph{residual-v3}; the specification is
\texttt{docs/RESIDUAL.md} in the artifact.

The normalizer is not incidental, and two earlier choices were discarded
before arriving at it. Normalizing by the signed grand sum
$\max(\lvert \mathbf{1}^{\!\top} C \mathbf{1} \rvert,
\lvert \mathbf{1}^{\!\top} A (B\mathbf{1}) \rvert)$ produces a heavy-tailed
quotient from two light-tailed pieces, because the denominator measures how
completely random signs cancelled rather than how large the multiply-adds
were. On our data that normalizer gave a max-to-median ratio of $389$, with
Pearson correlation $0.794$ between the residual and the reciprocal of its own
denominator. The tail was the normalizer, not the hardware.

Replacing the denominator with $e^{\!\top} \lvert A \rvert \lvert B \rvert e$
removes that failure mode. Every term is nonnegative, so cancellation cannot
drive it toward zero, and it concentrates tightly: across $n = 2000$ draws it
spanned $1.7169 \times 10^{10}$ to $1.7193 \times 10^{10}$, a max-to-median
ratio of $1.0007$. It is also the quantity the componentwise theory controls.
Writing $C = AB + E$, the checksum functional satisfies
$\lvert \mathbf{1}^{\!\top} E \mathbf{1} \rvert \leq
\gamma_K \, e^{\!\top} \lvert A \rvert \lvert B \rvert e$ under the classical
bound $\lvert E \rvert \leq \gamma_K \lvert A \rvert \lvert B \rvert$
\citep{higham2002accuracy}. We use this for its structural content only.
At bfloat16 with $K = 4096$ the product $Ku$ exceeds unity and $\gamma_K$ is
not a usable number, and tensor-core paths accumulate in a wider type before
storing, so the constant in front of $\lvert A \rvert \lvert B \rvert$ is not
the classical one either. The claim we rely on is that the checksum functional
is bounded by a nonnegative contraction of the factors, and that contraction
is the denominator we chose. Every numerical cutoff in this work comes from
measurement, not from $\gamma_K$.

The final change was to stop reducing before comparing. An earlier version of
this work summed the length-$M$ checksum vector to a scalar and compared
scalars, discarding $M-1$ degrees of freedom. Under that formulation, sign
flips and mantissa flips were undetectable in every cell we ran, $0$ of $1800$
samples, and we described the effect as structural. It was not structural. It
was the reduction. Under Equation~\ref{eq:residual-v3}, single sign flips are
caught in $155$ of $200$ samples on the same workload, the same shape, the
same injection harness and the same seeds. That earlier finding is retracted;
Section~\ref{sec:results:retractions} treats it and three others in full.

Two properties of Equation~\ref{eq:residual-v3} are worth stating because they
bound what the detector can be blamed for. The denominator does not depend on
$C$, so a fault in the product cannot shrink the scale that the fault is
measured against. And the per-row comparison is only informative if row
residuals are not dominated by a shared error term; on the $2000 \times 256$
matrix of per-row residuals we measure a mean pairwise Pearson correlation of
$0.000126$ (standard deviation $0.0225$, range $-0.0853$ to $0.0980$), so the
rows are effectively independent and the sensitivity gain over the scalar form
is not an artifact of correlated noise.

\subsection{Perturbation model}
\label{sec:method:perturbation}

Every detection number in this paper was produced under a single, versioned
perturbation model, specified in \texttt{docs/SPEC-PERTURBATION-MODEL.md} as
\texttt{perturbation-v1}. We state it early and plainly because it is the
strongest constraint on how our results may be read.

After $C = AB$ completes, distinct $(\text{element}, \text{bit})$ pairs inside
$C$ are XOR-flipped. The residual is then computed from the original factors
$A$ and $B$ and the perturbed product. Bits are drawn from one of five classes
of the bfloat16 layout, and the number of flips per sample is $1$, $2$, or $4$:

\begin{center}
\begin{tabular}{llc}
\toprule
Class & Bits & Width \\
\midrule
\texttt{SIGN}           & $[15]$    & 1 \\
\texttt{EXPONENT\_HIGH} & $[14{:}11]$ & 4 \\
\texttt{EXPONENT\_LOW}  & $[10{:}7]$  & 4 \\
\texttt{MANTISSA\_HIGH} & $[6{:}3]$   & 4 \\
\texttt{MANTISSA\_LOW}  & $[2{:}0]$   & 3 \\
\bottomrule
\end{tabular}
\end{center}

Nothing else is perturbed. Operand $A$ is never flipped, nor is operand $B$,
nor accumulator intermediates inside the GEMM kernel, nor instruction
encodings, nor control flow, nor any memory outside the output tensor. This is
an exhaustive list, and it has a direct consequence: there is no operand
asymmetry anywhere in our data, because every recorded flip perturbs the
numerator of Equation~\ref{eq:residual-v3} and leaves the denominator
untouched.

The obvious objection is that output-tensor XOR is not what a soft error does
to a GPU. Instruction-level and micro-architectural effects are invisible to an
output-tensor XOR model, and finer-grained injection studies operate below this
abstraction \citep{chai2026llm}. We concede the point
without qualification. This paper characterizes a detector under a defined
perturbation model. It does not characterize faults in silicon, and no rate
reported here estimates the incidence of anything in deployed hardware. The
fleet silent-corruption literature \citep{hochschild2021cores,dixit2021silent}
motivates why a cheap detector would be worth having; it does not supply our
numbers, and our numbers do not speak to it.

We verified that the injected perturbations are not silently discarded before
they can be measured. Mantissa flips are the case at risk, since a
low-mantissa flip followed by a cast could in principle round away. On
$n = 50$ samples per class, zero targeted elements were bitwise identical after
the cast, and zero samples had a maximum relative delta of zero. The
\texttt{MANTISSA\_HIGH} control produced relative deltas from $0.0317$ to
$0.492$ and \texttt{MANTISSA\_LOW} from $0.00407$ to $0.0313$: smaller, as
expected for lower bits, but strictly positive on every sample. The
undetectability reported in Section~\ref{sec:results:mantissa} is a real miss
of a real perturbation.

\subsection{Workload, hardware, and platform determinism}
\label{sec:method:platform}

All measurements use workload W02 in bfloat16 on a single Tesla T4 (sm75),
with \texttt{torch} 2.13.0+cu130, CUDA 13.0, and driver 580.159.04. Factors are
drawn as independent uniform $[-1, 1]$ matrices from a deterministic seed mixer
over base seed, workload identifier, shape index, and sample index, with no
wall-clock and no hashing, so every draw in this paper is reproducible from its
index alone.

Two shape regimes appear. The K-scaling study fixes $M = N = 4096$ and varies
$K \in \{512, 1024, 2048, 4096, 8192\}$, with $n = 500$ clean samples per $K$
and $n = 100$ flip samples per $(K, \text{class}, n_{\text{flips}})$ cell. The
threshold and false-positive analysis fixes $M = N = K = 4096$, with $n = 2000$
clean samples for the pilot floor, $n = 200$ flip samples per cell, and a
separate clean sweep of $20{,}000$ GEMMs described in
Section~\ref{sec:method:threshold}.

A detector whose clean floor is contaminated by library nondeterminism cannot
be characterized at all, so we checked. Across $44$ suite cases spanning every
GEMM shape in fp32, bf16 and fp16, attention, reductions, transcendentals and a
transformer block, all $44$ reported bitwise-identical output across repeated
launches on this machine. On this hardware, cuBLAS selected a stable kernel
every time, which supports treating the clean floor as accumulation error
rather than kernel-selection variance. We state the scope of that observation
precisely: it is a measured property of one tuple of GPU model, driver, CUDA
version and framework build. It is not a general property of cuBLAS and must
not be carried onto other hardware.

\subsection{Threshold and detection predicate}
\label{sec:method:threshold}

Residual-v3 performs $M = 4096$ per-row tests per GEMM. A per-GEMM
false-positive rate below $10^{-6}$, which is the bar this detector was built
against, therefore requires a per-row rate of
$1 - (1 - 10^{-6})^{1/4096} \approx 2.44 \times 10^{-10}$ under independence.
No sampling budget we could afford estimates that directly, so it has to be
extrapolated, and the extrapolation method was fixed in writing before the tail
was examined.

We collect clean per-row residuals from $20{,}000$ GEMMs at $K = 4096$, using
the full 4096-row vector rather than a subsample, for roughly
$8 \times 10^{7}$ row observations. Thresholds at the 95th, 99th and 99.9th
percentiles are frozen from an initial pass of 500 GEMMs, giving
$1.057096 \times 10^{-6}$, $1.392735 \times 10^{-6}$ and
$1.776693 \times 10^{-6}$. A Generalized Pareto distribution is fitted to
exceedances above each frozen threshold by maximum likelihood, chosen because
it is the standard peaks-over-threshold model and not because it fits these
data. The $1 - 2.44 \times 10^{-10}$ quantile of the fitted distribution is the
threshold. Fitting at three separate percentiles is a pre-registered
sensitivity check, and the pre-registered rule was that thresholds differing by
more than a factor of two are to be reported as untrustworthy rather than
averaged.

The canonical threshold is the extrapolation from the 99th-percentile fit,
%
\begin{equation}
\tau = 3.144535556 \times 10^{-6},
\end{equation}
%
and every detection rate reported in this paper is scored against it unless
explicitly stated otherwise. Section~\ref{sec:results:fpr} reports the fitted
parameters, the sensitivity comparison, and the one threshold that turned out
not to be evaluable.

Two points of provenance need stating here rather than in a limitations
section. First, $\tau$ is not the largest clean residual we observed. Observed
maxima are sample maxima, they grow with $n$, and two of them appear in our own
records for the same configuration at different sample sizes:
$3.245921 \times 10^{-6}$ at $n = 2000$ and $2.684525 \times 10^{-6}$ at
$n = 500$. Detection rates differ by a few points depending on which is used,
with no change to any rule or fit. We report $\tau$ as canonical because it is
the only threshold with a stated per-GEMM false-positive rate behind it, and we
note that it falls between the two sample maxima, so it is not an artifact of
either sampling run. Second, the conversion from per-GEMM to per-row rate
assumes independence across rows. Rather than assert it, we test it: the number
of per-row exceedances above the 99th percentile has observed mean $40.42$ and
variance $39.93$ per GEMM, against $40.96$ and $40.55$ for
$\mathrm{Binomial}(4096, 0.01)$, a variance ratio of $0.985$. The counts are
very slightly under-dispersed relative to independent Bernoulli trials. This
supports the conversion at the fitting threshold, and it does not establish
independence deep in the tail, which we carry forward as a stated assumption.

The detection predicate itself requires one sentence that would otherwise seem
too obvious to write. A sample is detected when $r_{\max} > \tau$
\emph{or} $r_{\max}$ is not finite. An infinite residual means the checksum
diverged, which is the loudest signal the detector can emit, and a naive
comparison against a threshold silently evaluates false on it. We found this in
our own analysis code during a pre-publication audit, where it had understated
detection in ten of fifteen \texttt{EXPONENT\_HIGH} cells by as much as
fourteen percentage points and manufactured an apparent effect that does not
exist. Section~\ref{sec:results:retractions} records the correction. We note
the general form of the hazard here because it applies to any threshold
detector operating on floating-point residuals: the case that most obviously
indicates failure is the case that comparison operators handle worst.

\subsection{Pre-registration}
\label{sec:method:prereg}

Every headline claim in this paper was predicted before it was measured.
Predictions were committed to the repository, with numeric ranges and an
explicit falsifier, before the corresponding data existed, and the commit
history establishes the ordering independently of our word for it. Four
prediction sets appear: the noise-floor and sensitivity predictions for
residual-v3, the K-scaling predictions, the extreme-value method and its
predictions, and a later mechanism prediction about operand asymmetry.

The rule we imposed on ourselves is that a prediction may be falsified by new
data but may not be rewritten by it. That rule was tested repeatedly and is the
reason this paper is worth reading rather than the reason it is short. Of the
K-scaling predictions, three of five failed, including the one that motivated
the experiment. The pre-registered falsifier for the underlying mechanism
fired. In the residual-v3 set, two of five predictions missed, one of them by
two percentage points, and it is reported as a miss rather than rounded to a
pass. A mechanism prediction registered late in the project was falsified at
its precondition when we discovered the harness could not test it, and it is
reported as such rather than deleted. The kill-test bar that this detector
failed is the same number that appeared in the repository README before any
code was written, and it has never been moved.

We also state where pre-registration did not save us. Two defects reached
published internal documents and were caught by audit rather than by
prediction. The first is the detection predicate described above, withdrawn in
Section~\ref{sec:results:retractions}. The second was a results document citing
a source data file under a path that never existed; the file was located, the
numbers it produced were reproduced from it, and the provenance was corrected
in the repository. Writing the bar down first constrains
what you may conclude from a measurement. It does not guarantee the measurement
is right.
